\documentclass[12pt,a4paper]{article}

\usepackage[T1]{fontenc}
\usepackage[utf8]{inputenc}
\usepackage[brazil]{babel}
\usepackage{lmodern}
\usepackage{geometry}
\usepackage{graphicx}
\usepackage{booktabs}
\usepackage{array}
\usepackage{xcolor}
\usepackage{hyperref}
\usepackage{float}
\usepackage{enumitem}
\usepackage{fancyhdr}

\geometry{top=2.5cm,bottom=2.5cm,left=2.5cm,right=2.5cm}
\graphicspath{{TestResults/SeleniumScreenshots/}}
\hypersetup{
  colorlinks=true,
  linkcolor=blue!50!black,
  urlcolor=blue!60!black,
  pdftitle={Relatório de Testes Selenium --- BiotLab},
  pdfauthor={Equipe BiotLab}
}

\definecolor{biotblue}{HTML}{2563EB}
\definecolor{successgreen}{HTML}{15803D}
\definecolor{lightgray}{HTML}{F3F4F6}

\pagestyle{fancy}
\fancyhf{}
\lhead{\textbf{BiotLab}}
\rhead{Relatório de Testes Selenium}
\cfoot{\thepage}
\setlength{\headheight}{15pt}

\newcommand{\aprovado}{\textcolor{successgreen}{\textbf{Aprovado}}}

\begin{document}

\begin{titlepage}
  \centering
  \vspace*{2.5cm}
  {\color{biotblue}\rule{\textwidth}{2pt}\par}
  \vspace{1.2cm}
  {\Huge\bfseries Relatório de Testes Selenium\par}
  \vspace{0.5cm}
  {\LARGE Sistema BiotLab\par}
  \vspace{1.5cm}
  {\large Testes funcionais automatizados de interface\par}
  \vfill
  \begin{tabular}{rl}
    \textbf{Data da execução:} & 17 de julho de 2026 \\
    \textbf{Aplicação:} & BiotLab Web \\
    \textbf{Tecnologia:} & ASP.NET Core .NET 8 \\
    \textbf{Ferramenta:} & Selenium WebDriver 4.46.0 \\
    \textbf{Navegador:} & Google Chrome, modo headless \\
    \textbf{Banco de dados:} & MySQL local \\
  \end{tabular}
  \vfill
  {\color{biotblue}\rule{\textwidth}{2pt}\par}
\end{titlepage}

\tableofcontents
\newpage

\section{Resumo executivo}

Foram executados três testes funcionais automatizados sobre a interface web do
BiotLab. Todos os cenários foram aprovados, incluindo a apresentação da tela de
login, a proteção de uma rota contra acesso anônimo e a autenticação de um
administrador com acesso posterior ao dashboard.

\begin{center}
  \setlength{\tabcolsep}{12pt}
  \renewcommand{\arraystretch}{1.4}
  \begin{tabular}{>{\bfseries}l c}
    \toprule
    Indicador & Resultado \\
    \midrule
    Testes executados & 3 \\
    Testes aprovados & \textcolor{successgreen}{\textbf{3}} \\
    Testes com falha & 0 \\
    Testes ignorados & 0 \\
    Duração aproximada & 13 segundos \\
    \bottomrule
  \end{tabular}
\end{center}

\section{Objetivo e escopo}

O objetivo desta execução foi validar os fluxos essenciais de acesso ao sistema
por meio de um navegador real automatizado. O escopo contemplou:

\begin{enumerate}[label=\arabic*.]
  \item carregamento da página de autenticação;
  \item presença dos campos e controles obrigatórios;
  \item bloqueio de acesso anônimo a um módulo protegido;
  \item redirecionamento para a página de login;
  \item autenticação administrativa;
  \item apresentação do dashboard após o login.
\end{enumerate}

\section{Resultados}

\renewcommand{\arraystretch}{1.35}
\begin{table}[H]
  \centering
  \small
  \begin{tabular}{p{4.2cm} p{2.2cm} p{7.2cm}}
    \toprule
    \textbf{Cenário} & \textbf{Resultado} & \textbf{Evidência validada} \\
    \midrule
    Carregamento da página de login
      & \aprovado
      & Formulário, campo de e-mail, campo de senha e botão de entrada visíveis e habilitados. \\
    Proteção de módulo para usuário anônimo
      & \aprovado
      & O acesso a \texttt{/Instituicao} foi redirecionado ao login com parâmetro \texttt{ReturnUrl}. \\
    Login administrativo
      & \aprovado
      & As credenciais foram aceitas, a página de login foi encerrada e o dashboard do BiotLab foi exibido. \\
    \bottomrule
  \end{tabular}
  \caption{Resumo dos cenários executados pelo Selenium.}
\end{table}

O resultado estruturado da execução também foi armazenado no arquivo
\path{BiotLabTests/TestResults/selenium-results.trx}.

\section{Evidências visuais}

\subsection{Página de login}

A Figura~\ref{fig:login} demonstra a apresentação da página de autenticação
com os controles exigidos pelo cenário de teste.

\begin{figure}[H]
  \centering
  \includegraphics[width=\textwidth]{01-pagina-login.png}
  \caption{Página de login carregada pelo Selenium.}
  \label{fig:login}
\end{figure}

\subsection{Redirecionamento de acesso anônimo}

A Figura~\ref{fig:redirect} registra o resultado da tentativa de acesso anônimo
a um módulo protegido. O sistema encaminhou o navegador para a autenticação.

\begin{figure}[H]
  \centering
  \includegraphics[width=\textwidth]{02-redirecionamento-login.png}
  \caption{Redirecionamento de usuário anônimo para o login.}
  \label{fig:redirect}
\end{figure}

\subsection{Dashboard administrativo}

A Figura~\ref{fig:dashboard} comprova que o administrador autenticado chegou
ao dashboard e recebeu os módulos correspondentes ao seu perfil.

\begin{figure}[H]
  \centering
  \includegraphics[width=\textwidth]{03-dashboard-administrador.png}
  \caption{Dashboard apresentado após autenticação administrativa.}
  \label{fig:dashboard}
\end{figure}

\section{Avaliação}

Os resultados demonstram que os fluxos essenciais cobertos estão funcionando:

\begin{itemize}
  \item a página de autenticação oferece os controles necessários;
  \item rotas protegidas impedem o acesso de usuários anônimos;
  \item um administrador válido consegue autenticar-se;
  \item o sistema apresenta o dashboard após a autenticação;
  \item os testes podem produzir evidências visuais e um resultado TRX.
\end{itemize}

\section{Limitações}

Esta execução representa uma verificação inicial de fumaça e ainda não cobre:

\begin{itemize}
  \item logout, expiração e reutilização de sessão;
  \item bloqueio após repetidas tentativas inválidas;
  \item recuperação e redefinição de senha;
  \item permissões específicas do perfil de estudante;
  \item fluxos de cadastro, consulta, alteração e exclusão;
  \item mensagens e validações dos formulários;
  \item navegadores, dispositivos e resoluções adicionais;
  \item acessibilidade, desempenho e carga.
\end{itemize}

\section{Risco de segurança identificado}

Durante a preparação do ambiente foram observadas credenciais armazenadas
diretamente em arquivos \texttt{appsettings}. Nenhuma credencial é reproduzida
neste relatório.

Recomenda-se:

\begin{enumerate}[label=\arabic*.]
  \item revogar ou rotacionar imediatamente as credenciais expostas;
  \item remover senhas dos arquivos versionados;
  \item utilizar variáveis de ambiente ou .NET User Secrets;
  \item empregar uma conta exclusiva e de privilégio mínimo para automação;
  \item impedir que capturas, logs e relatórios registrem informações sensíveis.
\end{enumerate}

\section{Próximos testes recomendados}

\begin{enumerate}[label=\arabic*.]
  \item validar login inválido e bloqueio de conta;
  \item validar logout e impossibilidade de reutilizar a sessão;
  \item verificar autorização para cada perfil;
  \item automatizar os fluxos CRUD dos módulos administrativos;
  \item capturar automaticamente telas e HTML quando um cenário falhar;
  \item integrar os testes a uma rotina de integração contínua;
  \item acompanhar os resultados por execução usando TRX e relatório HTML.
\end{enumerate}

\section{Conclusão}

Os três cenários executados foram aprovados, sem falhas ou testes ignorados.
Dentro do escopo avaliado, o BiotLab apresentou corretamente a interface de
login, o controle de acesso anônimo e o dashboard administrativo. A ampliação
da suíte para operações de negócio e controles de segurança permitirá uma
avaliação funcional mais abrangente.

\end{document}
